%% kapa: exam=2 round=36 year=2025 period=3 subject=compensation_law subject_ko=감정평가 및 보상법규 seq=3 number=3 points=20
%% source: Q-Net 감정평가사 기출문제 (cst00309) artlSeq=5239325 file=raw/1008/5239325/2025년 제36회 감정평가사 2차 감정평가 및 보상법규.pdf
%% conversion: kordoc 4.24.51 -> tools/gen2.py (auto); review=proofread-gpt-sol; flags=none
\begin{longquestion}{【문제 3】}
국토교통부장관은 ｢감정평가 및 감정평가사에 관한 법률｣ 제6조(감정평가서)에 반하는 감정평가서를 작성ㆍ발급하였다는 이유로 감정평가사 甲에게 3월의 업무정지처분을 하였다. 甲은 위 처분에 대하여 행정심판을 청구하였고, 행정심판위원회는 “국토교통부장관은 甲에게 한 3월의 업무정지처분을 2월의 업무정지처분으로 변경하라”는 재결을 하였다. 국토교통부장관은 위 재결취지에 따라 “3월의 업무정지처분을 2월의 업무정지처분으로 변경한다”는 후속 변경처분을 하였다. 다음 물음에 답하시오. (20점) \par
\medskip\noindent 물음 1) 甲이 후속 변경처분을 받은 후에 취소소송을 제기하는 경우 어떠한 처분을 대상으로 하는지 검토하시오. (10점) \par
\medskip\noindent 물음 2) 위 행정쟁송이 있기 전, 국토교통부장관은 ｢감정평가 및 감정평가사에 관한 법률｣상 업무정지처분을 갈음하여 과징금을 부과할 수 있는지, 그리고 甲이 과징금 납부의무를 불이행한 경우 국토교통부장관은 어떠한 조치를 할 수 있는지 검토하시오. (10점) \par
\end{longquestion}
